\chapter*{Abstract}
\addcontentsline{toc}{chapter}{Abstract}

The nonuniform fast Fourier transform underlies image reconstruction in magnetic
resonance imaging, gridding in radio interferometry and particle-mesh methods in
molecular dynamics, and FINUFFT is among the fastest implementations of it. How
the library computes a transform has been studied closely; how it should be
configured has not. The library exposes a set of options and
fixes several further constants in its source; six of them decide how the work is
divided and how accurately each piece is computed. Their values come from
constants chosen once, by hand, for a machine and a workload that may not be the
ones in front of it.

This thesis measures what those parameters do and why. The library is
instrumented so that the cost of each stage of a transform can be attributed to
that step alone, using timers placed inside its parallel region together with
hardware performance counters. A search over the joint parameter space
establishes which parameters account for the variation in execution time, and
each of those is then examined on its own across dimensions, point densities,
requested tolerances and thread counts. Where a cost follows from the structure
of the algorithm, it is derived in closed form and checked against measurement.

The parameters are found to act in opposition rather than in isolation: each
makes one step of the transform cheaper only by making another more expensive,
so every optimum lies between the extremes rather than at either end, and its
position moves with the problem and with the machine. The distances involved are
large. One parameter is set fifty times above its best value for the transform
measured, which costs a factor of $1.44$ in the phase it governs; a search over
six parameters together finds a configuration $19\%$ faster than the defaults.
The cause is usually the memory hierarchy: a stage doubles in cost at the point
where the data a thread holds outgrows its share of the second-level cache, and
a second parameter, worth nothing at the library's default, decides up to a
third of the run time once the first is tuned.

The behaviour therefore divides in two. One part follows from the geometry of
the grids the algorithm builds, and is predictable: a closed-form model of the
subgrid a configuration produces matches measurement to within a percent, and
with it the cost of a step that was previously only measurable. The other part
follows from where the data sits in the cache hierarchy; its thresholds can be
located from the cache sizes, but the cost of crossing them must be measured on
the machine at hand. The contribution is an account of why each parameter
behaves as it does, and of the quantities on which an automatic choice would
have to be based.

\vspace{1cm}
\noindent\textbf{Keywords:} FINUFFT, nonuniform FFT, parameter selection, performance analysis
